% Préambule commun des schémas TikZ — charte Liouaï
% Chaque schéma commence par % Préambule commun des schémas TikZ — charte Liouaï
% Chaque schéma commence par % Préambule commun des schémas TikZ — charte Liouaï
% Chaque schéma commence par \input{_preambule} puis \begin{document}.
\documentclass[tikz,border=6pt]{standalone}
\usepackage[T1]{fontenc}
\usepackage[utf8]{inputenc}
\usepackage[scaled=0.95]{helvet}
\renewcommand{\familydefault}{\sfdefault}
\usepackage{amsmath,amssymb}
\usepackage[european]{circuitikz}
\usetikzlibrary{arrows.meta,positioning,calc,angles,quotes,decorations.pathreplacing,backgrounds}

% Couleurs du logo Liouaï Digital Education
\definecolor{navy}{HTML}{243743}
\definecolor{vert}{HTML}{28B78D}
\definecolor{vertfonce}{HTML}{167A5C}
\definecolor{vertclair}{HTML}{EAF7F2}
\definecolor{gris}{HTML}{A1A9AE}
\definecolor{gristexte}{HTML}{5C6870}
\definecolor{fond}{HTML}{F3F5F6}
\definecolor{jaune}{HTML}{E5C800}
\definecolor{jauneclair}{HTML}{FFFBD6}

\tikzset{
  every picture/.style={line cap=round, line join=round, color=navy, text=navy},
  fleche/.style={-{Stealth[length=2.6mm,width=2mm]}, line width=1.1pt},
  vecteur/.style={fleche, line width=1.4pt},
  titre/.style={font=\bfseries\small, text=navy},
  note/.style={font=\footnotesize, text=gristexte},
  boite/.style={draw=navy, line width=0.9pt, rounded corners=5pt, fill=fond, align=center, inner sep=6pt},
}
 puis \begin{document}.
\documentclass[tikz,border=6pt]{standalone}
\usepackage[T1]{fontenc}
\usepackage[utf8]{inputenc}
\usepackage[scaled=0.95]{helvet}
\renewcommand{\familydefault}{\sfdefault}
\usepackage{amsmath,amssymb}
\usepackage[european]{circuitikz}
\usetikzlibrary{arrows.meta,positioning,calc,angles,quotes,decorations.pathreplacing,backgrounds}

% Couleurs du logo Liouaï Digital Education
\definecolor{navy}{HTML}{243743}
\definecolor{vert}{HTML}{28B78D}
\definecolor{vertfonce}{HTML}{167A5C}
\definecolor{vertclair}{HTML}{EAF7F2}
\definecolor{gris}{HTML}{A1A9AE}
\definecolor{gristexte}{HTML}{5C6870}
\definecolor{fond}{HTML}{F3F5F6}
\definecolor{jaune}{HTML}{E5C800}
\definecolor{jauneclair}{HTML}{FFFBD6}

\tikzset{
  every picture/.style={line cap=round, line join=round, color=navy, text=navy},
  fleche/.style={-{Stealth[length=2.6mm,width=2mm]}, line width=1.1pt},
  vecteur/.style={fleche, line width=1.4pt},
  titre/.style={font=\bfseries\small, text=navy},
  note/.style={font=\footnotesize, text=gristexte},
  boite/.style={draw=navy, line width=0.9pt, rounded corners=5pt, fill=fond, align=center, inner sep=6pt},
}
 puis \begin{document}.
\documentclass[tikz,border=6pt]{standalone}
\usepackage[T1]{fontenc}
\usepackage[utf8]{inputenc}
\usepackage[scaled=0.95]{helvet}
\renewcommand{\familydefault}{\sfdefault}
\usepackage{amsmath,amssymb}
\usepackage[european]{circuitikz}
\usetikzlibrary{arrows.meta,positioning,calc,angles,quotes,decorations.pathreplacing,backgrounds}

% Couleurs du logo Liouaï Digital Education
\definecolor{navy}{HTML}{243743}
\definecolor{vert}{HTML}{28B78D}
\definecolor{vertfonce}{HTML}{167A5C}
\definecolor{vertclair}{HTML}{EAF7F2}
\definecolor{gris}{HTML}{A1A9AE}
\definecolor{gristexte}{HTML}{5C6870}
\definecolor{fond}{HTML}{F3F5F6}
\definecolor{jaune}{HTML}{E5C800}
\definecolor{jauneclair}{HTML}{FFFBD6}

\tikzset{
  every picture/.style={line cap=round, line join=round, color=navy, text=navy},
  fleche/.style={-{Stealth[length=2.6mm,width=2mm]}, line width=1.1pt},
  vecteur/.style={fleche, line width=1.4pt},
  titre/.style={font=\bfseries\small, text=navy},
  note/.style={font=\footnotesize, text=gristexte},
  boite/.style={draw=navy, line width=0.9pt, rounded corners=5pt, fill=fond, align=center, inner sep=6pt},
}

\begin{document}
\begin{tikzpicture}[
  secteur/.style={circle, draw=navy, line width=1pt, fill=vertclair, minimum size=2.1cm, align=center, font=\small},
  flux/.style={fleche, vertfonce, line width=1pt},
  coef/.style={font=\footnotesize, fill=white, inner sep=1.5pt}]
\node[secteur] (ag) at (0,0) {Agriculture};
\node[secteur] (in) at (5,0) {Industrie};
\node[secteur] (se) at (2.5,-4) {Services};
% fournitures : du fournisseur vers le client (coefficient a_ij : ce que j achète à i pour produire 1 €)
\draw[flux] (ag) to[bend left=18] node[coef] {0,15} (in);
\draw[flux] (in) to[bend left=18] node[coef] {0,20} (ag);
\draw[flux] (in) to[bend left=18] node[coef] {0,10} (se);
\draw[flux] (se) to[bend left=18] node[coef] {0,20} (in);
\draw[flux] (se) to[bend left=18] node[coef] {0,10} (ag);
\draw[flux] (ag) to[bend left=18] node[coef] {0,02} (se);
% autoconsommation
\draw[flux] (ag) to[out=150, in=210, looseness=5] node[coef, left] {0,10} (ag);
\draw[flux] (in) to[out=30, in=-30, looseness=5] node[coef, right] {0,30} (in);
\draw[flux] (se) to[out=-120, in=-60, looseness=5] node[coef, left] {0,25} (se);
% demande finale : une flèche sortante par secteur
\draw[fleche, navy, line width=1.3pt] (ag) -- (-2.4,1.6) node[left, font=\footnotesize] {$d_1$};
\draw[fleche, navy, line width=1.3pt] (in) -- (7.4,1.6) node[right, font=\footnotesize] {$d_2$};
\draw[fleche, navy, line width=1.3pt] (se) -- (5.6,-4.9) node[right, font=\footnotesize] {$d_3$};
\node[note, align=left, anchor=west] at (-1.6,-7.4) {Flèches vertes : ce qu'un secteur fournit à un autre pour chaque euro produit\\ par le client (exemple : pour produire 1\,€, l'industrie achète 0,15\,€ à l'agriculture).\\ Flèches marine : la demande finale (ménages, exportations). Valeurs illustratives.};
\end{tikzpicture}
\end{document}
